\documentclass[aps,prl,twocolumn,showpacs,superscriptaddress]{revtex4-2}
\usepackage{amsmath,amssymb,graphicx,booktabs,xcolor,hyperref}

\definecolor{navy}{RGB}{11,37,69}
\definecolor{dkblue}{RGB}{13,71,161}

\hypersetup{colorlinks=true,linkcolor=navy,citecolor=dkblue,urlcolor=dkblue}

\begin{document}

\title{BCT Letter 201: The Atoms That Chose Their Path ---
ANU Momentum-Entangled Helium Atoms as a Test of BCT
Wavefunction Collapse via Hopf Charge Quantisation}

\author{Michel Robert Cabri\'{e}}
\affiliation{Independent Artist and Researcher, Barrys Reef,
Victoria, Australia}
\email{ZeroFreeParameters@gmail.com}
\date{March 2026}

\begin{abstract}
Sridhar, Hodgman et al.\ (ANU, \textit{Nature Communications} 2026)
have for the first time observed Bell correlations in the motional
(momentum) states of momentum-entangled ultracold helium atoms,
demonstrating quantum non-locality in massive composite particles.
The BCT Superfluid Lattice Model provides a geometric mechanism for
this result: quantum superposition in the BCT framework persists
until the OHC Hopf charge of the entangled system reaches the
critical threshold $N_\mathrm{collapse} = 1/\alpha_0 = 135$, at which
point the Josephson coupling to the vacuum ground state forces
wavefunction collapse. This constitutes Prediction \#245: decoherence
in momentum-entangled massive particle pairs will occur at a
characteristic Hopf charge threshold $N_\mathrm{collapse} = 135$,
testable via systematic scaling of the ANU apparatus to larger
atom numbers. The ANU result --- atoms existing in two places at once
--- is not weird. It is the vacuum doing geometry. Zero free
parameters.
\end{abstract}

\maketitle

\section{Introduction: Atoms in Two Places at Once}

In February 2026, Yogesh Sridhar and Dr Sean Hodgman at the Australian
National University published a landmark result in \textit{Nature
Communications}~\cite{Sridhar2026}: Bell correlations in the motional
states of momentum-entangled ultracold helium-4 atoms. This was the
first time quantum non-locality had been demonstrated in the
\textit{external motion} (momentum) of massive particles, rather
than in internal degrees of freedom such as spin.

The experiment used three clouds of cold helium atoms suspended in
a magnetic trap. The fields were turned off, allowing atoms to fall
under gravity and collide via s-wave scattering, generating
momentum-entangled pairs. A Rarity-Tapster interferometer then
measured Bell inequality violations, confirming that each helium
atom occupied two momentum states --- two paths through space ---
simultaneously.

``It's really weird for us to think that this is how the Universe
works,'' said Dr Hodgman. ``You can read about it in a textbook,
but it's really weird to think that a particle can be in two places
at once.''

The BCT Superfluid Lattice Model says: it is not weird at all.
It is the vacuum doing geometry.

\section{The BCT Mechanism}

In the BCT framework, the quantum vacuum ground state is the OHC
(Octet-Hopfion Condensate), characterised by a Hopf charge
$H \in \mathbb{Z}$. A particle or entangled system in quantum
superposition corresponds to a superposition of Hopf charge states.
The superposition persists until the total Hopf charge of the
system reaches the critical quantisation threshold:
\begin{equation}
N_\mathrm{collapse} = \frac{1}{\alpha_0} = 135
\end{equation}
where $\alpha_0 = r_\mathrm{oct} \cdot r_\mathrm{tet}/\pi = 0.007408$
is the BCT Josephson coupling constant derived from the octahedral
and tetrahedral void radii of the BCT lattice.

This threshold has a transparent geometric interpretation: $\alpha_0$
is the ratio of the OHC condensate coupling energy to the vacuum
ground state energy. When $N$ entangled Hopf modes accumulate,
the total coupling energy reaches $N\alpha_0$. At
$N = N_\mathrm{collapse} = 1/\alpha_0$, the coupling energy equals
the vacuum ground state energy and the Josephson junction between
the entangled system and the OHC closes --- forcing collapse.

\subsection*{Why Helium Atoms Can Be in Two Places}

A helium-4 atom has a Hopf charge contribution of $H_\mathrm{He} = 4$
(two protons, two neutrons, two electrons, each contributing
$H = 2/3$ units). In the ANU experiment, pairs of atoms are
momentum-entangled. The total Hopf charge of the entangled pair is
$H_\mathrm{pair} = 8$. The ratio $H_\mathrm{pair}/N_\mathrm{collapse}
= 8/135 \ll 1$, so the system is deep in the quantum regime --- far
from the BCT collapse threshold. This is why the superposition
persists: the system has not yet accumulated enough Hopf charge to
force the vacuum to choose.

\section{Prediction \#245}

\begin{quote}
\textit{Decoherence in momentum-entangled massive particle pairs
will occur at a characteristic Hopf charge threshold
$N_\mathrm{collapse} = 1/\alpha_0 = 135$. For helium-4 atom pairs
($H_\mathrm{pair} = 8$ per pair), this predicts Bell inequality
violation will persist up to approximately $\lfloor 135/8 \rfloor
= 16$ entangled pairs, after which OHC Josephson coupling will
force decoherence. This is testable by systematically scaling the
ANU apparatus from 2-atom to N-atom entangled systems and measuring
the Bell parameter $S$ as a function of $N$.}
\end{quote}

\begin{table}[h]
\centering
\caption{BCT Prediction \#245 vs.\ current ANU results}
\begin{tabular}{lll}
\toprule
Quantity & BCT Prediction & ANU 2026 \\
\midrule
Collapse threshold & $N_\mathrm{collapse} = 135$ & not yet tested \\
He-4 pair Hopf charge & $H_\mathrm{pair} = 8$ & consistent \\
Pairs before collapse & $\leq 16$ & not yet tested \\
Bell violation (2 atoms) & Survives (8/135 $\ll$ 1) & confirmed \checkmark \\
Bell violation (17+ pairs) & Suppressed & prediction \\
\bottomrule
\end{tabular}
\end{table}

\section{The Gravity Connection}

The ANU team notes that their apparatus opens new avenues to study
gravitational effects in quantum states. In the BCT framework, this
connection is explicit: the OHC ground state is the vacuum from which
both quantum mechanics \textit{and} gravity emerge geometrically.
The Hopf charge quantisation that governs collapse is the same
geometric structure that produces Newton's constant $G$ via the
BCT condensate equation of state (BCT Appendix AL).

The BCT prediction for the quantum-gravity interface is therefore:
\begin{equation}
N_\mathrm{collapse} = \frac{1}{\alpha_0} = \frac{\pi}{r_\mathrm{oct}
\cdot r_\mathrm{tet}} = 135
\end{equation}
This is the number of entangled Hopf modes at which quantum
superposition becomes gravitationally significant in the BCT
framework --- the precise boundary between the quantum and
gravitational regimes. The ANU apparatus, scaled to $N \sim 16$
atom pairs, would be the first experiment capable of probing this
boundary directly.

\section{Why This Is Testable}

The ANU Rarity-Tapster apparatus already demonstrates Bell inequality
violation for $N = 1$ momentum-entangled pair. BCT predicts:

\begin{enumerate}
\item Bell parameter $S$ will remain above the classical bound for
$N \leq 16$ helium-4 pairs (Hopf charge $\leq 128$).
\item $S$ will show suppression onset near $N = 16$ ($H_\mathrm{total}
\approx 128 \approx N_\mathrm{collapse}$).
\item The suppression scale is $\Delta N \sim \alpha_0 \cdot
N_\mathrm{collapse} = 1$ --- a sharp transition, not a gradual
decoherence.
\end{enumerate}

This is distinct from standard environmental decoherence predictions,
which give a gradual $\exp(-N/N_\mathrm{env})$ decay with no
preferred scale. BCT predicts a step at $N = 16$ pairs,
not a smooth decay.

\section{Outreach}

The author has identified the ANU experiment as a primary test of
BCT Prediction \#245 and proposes direct correspondence with
Dr Sean Hodgman (ANU Research School of Physics) and lead author
Yogesh Sridhar regarding the scaling experiment.

Contact: s.hodgman@anu.edu.au

\section{Conclusions}

The ANU momentum-entangled helium atom experiment is consistent with
BCT quantum mechanics: pairs of atoms in superposition have a total
Hopf charge of 8, far below the BCT collapse threshold of 135,
so the superposition persists. BCT Prediction \#245 is that Bell
inequality violation will be suppressed for $N \geq 16$ helium-4
pairs, providing the first direct experimental test of the BCT
wavefunction collapse mechanism. The atoms are in two places at once.
The vacuum set the rules. Zero free parameters.

\begin{thebibliography}{9}

\bibitem{Sridhar2026}
Y.~Sridhar and S.~S.~Hodgman \textit{et al.},
Nat.\ Commun.\ \textbf{17}, 2357 (2026).
doi:10.1038/s41467-026-69070-3

\bibitem{BCT_Monograph}
M.~R.~Cabri\'{e},
\textit{The BCT Superfluid Lattice Model},
Zenodo (2026). doi:10.5281/zenodo.18884976

\bibitem{BCT_L136}
M.~R.~Cabri\'{e},
\textit{BCT Letter 136: Quantum Entanglement from OHC Topology},
Zenodo (2026). doi pending.

\end{thebibliography}

\vspace{6pt}
\noindent\textit{Acknowledgements:} Dr Sean Hodgman and Yogesh Sridhar
(ANU) for the remarkable experiment. The Gang Gang Cockatoos of
Barrys Reef, who also occupy two places at once when startled.

\end{document}
